\documentclass{fiche}
\metadonnees{19}{La fenêtre}{Bloc 4 — Hypothèse, parole rapportée, traduction}

\begin{document}
\entete

\objectifs{%
\textbf{Langue} : les tours que le français rend mal --- \emph{on}, \emph{il y a},
\emph{depuis}, \emph{c'est\ldots qui} ; place de l'adverbe ; faux amis.\par
\textbf{Texte} : James Joyce, \og Eveline \fg, \emph{Dubliners} (1914).\par
\textbf{Durée} : 1\,h\,30 — exercices 35 min, texte et questions 30 min, version 25 min.}

%% ===================================================================
\exercice[14 min]{Quatre tours français qui n'ont pas d'équivalent}

\rappel{\textbf{Rappel.} \og On \fg{} se rend par un passif, par \emph{you} ou \emph{one},
ou par \emph{people} ; \og il y a \fg{} par \emph{there is} ou par \emph{ago} ; \og
depuis \fg{} par \emph{for} ou \emph{since} avec un present perfect ; \og c'est\ldots qui \fg{}
par l'accent, non par une construction.}

\consigne{Traduisez en anglais.}

\begin{anglais}
\begin{enumerate}[label=\arabic*.,itemsep=2.5mm,leftmargin=*]
  \item On lui a dit de rentrer.
    \lignes{1}
    \corr{She was told to go home.}
  \item On ne sait jamais.
    \lignes{1}
    \corr{You never know. \emph{ou} One never knows.}
  \item Ici on parle anglais.
    \lignes{1}
    \corr{English is spoken here.}
  \item Il y a une lampe sur la table.
    \lignes{1}
    \corr{There is a lamp on the table.}
  \item Il est parti il y a trois ans.
    \lignes{1}
    \corr{He left three years ago.}
  \item Elle habite là depuis dix-neuf ans.
    \lignes{1}
    \corr{She has lived there for nineteen years.}
  \item Elle attendait depuis une heure quand il arriva.
    \lignes{1}
    \corr{She had been waiting for an hour when he arrived.}
  \item C'est son père qui a acheté le champ.
    \lignes{1}
    \corr{It was her father who bought the field. \emph{ou} Her father was the one who bought
    the field.}
  \item C'est demain qu'elle part.
    \lignes{1}
    \corr{She is leaving tomorrow. \emph{(l'accent porte sur} tomorrow\emph{)}}
  \item On dit qu'il est à Melbourne.
    \lignes{1}
    \corr{He is said to be in Melbourne. \emph{(fiche 14)}}
\end{enumerate}
\end{anglais}

\note[Le piège de \og depuis \fg]{Le français emploie le présent (\og elle habite là
depuis\ldots \fg), l'anglais le present perfect, parce qu'il regarde la période écoulée et
non l'instant présent. Au passé, le même décalage : \og elle attendait depuis une heure \fg{}
$\rightarrow$ \emph{she had been waiting for an hour}. C'est la faute la plus coûteuse du
thème.}

\note[Le piège de \og c'est\ldots qui \fg]{Le français met en relief par une construction,
l'anglais par l'accent de phrase et par l'ordre des mots. \emph{It was her father who\ldots}
existe, mais l'anglais le réserve à une vraie opposition (et non un autre). Le plus souvent
il se contente de placer le mot important à la fin, ou de l'appuyer à l'oral.}

%% ===================================================================
\exercice[10 min]{Place de l'adverbe}

\rappel{\textbf{Rappel.} Les adverbes de fréquence se placent avant le verbe, mais après
\emph{be} et après l'auxiliaire. Rien ne se glisse jamais entre le verbe et son complément
d'objet.}

\consigne{Remettez l'adverbe à sa place.}

\begin{anglais}
\begin{enumerate}[label=\arabic*.,itemsep=2.2mm,leftmargin=*]
  \item She dusted the room. \emph{(once a week)}
    \trou{She dusted the room once a week.}
  \item Her father used to hunt them out of the field. \emph{(often)}
    \trou{Her father used often to hunt them out of the field.}
  \item He was patient with her. \emph{(always)}
    \trou{He was always patient with her.}
  \item She had found out his name. \emph{(never)}
    \trou{She had never found out his name.}
  \item She played with the other children. \emph{(every evening)}
    \trou{She played with the other children every evening.}
  \item He goes for her now. \emph{(seldom)}
    \trou{He seldom goes for her now.}
  \item She would cry many tears. \emph{(not)}
    \trou{She would not cry many tears.}
  \item They treated her with respect. \emph{(never)}
    \trou{They never treated her with respect.}
\end{enumerate}
\end{anglais}

\note[Trois règles]{Adverbe de fréquence : avant le verbe lexical (\emph{she never cries}),
après \emph{be} (\emph{he is always late}), après le premier auxiliaire (\emph{she had never
found out}). Adverbe de temps précis : en fin de phrase (\emph{once a week}, \emph{every
evening}). Et jamais entre le verbe et son objet : \emph{she dusted carefully the room} est
impossible.}

%% ===================================================================
\exercice[11 min]{Faux amis}
\consigne{Traduisez, puis donnez le mot anglais correspondant au faux sens.}

\begin{center}\small
\begin{tabular}{@{}lll@{}}
\textsf{mot anglais} & \textsf{sens réel} & \textsf{et \og\,\ldots\,\fg{} se dit} \\[1.5mm]
\emph{actually} & \trou{en fait} & \trou{currently} \\[1mm]
\emph{eventually} & \trou{finalement} & \trou{possibly} \\[1mm]
\emph{to demand} & \trou{exiger} & \trou{to ask} \\[1mm]
\emph{sensible} & \trou{raisonnable} & \trou{sensitive} \\[1mm]
\emph{a library} & \trou{une bibliothèque} & \trou{a bookshop} \\[1mm]
\emph{to support} & \trou{soutenir} & \trou{to bear, to stand} \\[1mm]
\emph{a deception} & \trou{une tromperie} & \trou{a disappointment} \\[1mm]
\emph{an issue} & \trou{une question, une sortie} & \trou{an outcome} \\[1mm]
\emph{to pretend} & \trou{feindre} & \trou{to claim} \\[1mm]
\emph{a figure} & \trou{un chiffre, une silhouette} & \trou{a face} \\[1mm]
\emph{to achieve} & \trou{accomplir} & \trou{to finish} \\[1mm]
\emph{a college} & \trou{un établissement supérieur} & \trou{a secondary school} \\
\end{tabular}
\end{center}

\note{Trois de ces mots ont déjà servi : \emph{eventually} en fiche 5, \emph{to demand} en
fiche 12, \emph{a figure} en fiche 11. Un faux ami ne se corrige qu'à force de le rencontrer
juste.}

%% ===================================================================
\newpage
\section{Texte}

\noindent\small\textit{Eveline Hill, dix-neuf ans, vit à Dublin avec son père. Elle doit
partir le soir même pour Buenos Aires avec Frank, un marin qu'elle veut épouser.}\normalsize

\begin{texte}
She sat at the window watching the evening invade the avenue. Her head was leaned against the
window curtains and in her nostrils was the odour of dusty
\gl[cretonne]{cretonne}{la cretonne, tissu d'ameublement}. She was tired.

Few people passed. The man out of the last house passed on his way home; she heard his
footsteps clacking along the concrete pavement and afterwards
\gl[to crunch]{crunching}{crisser} on the \gl[a cinder path]{cinder path}{une allée de
mâchefer} before the new red houses. One time there used to be a field there in which they
used to play every evening with other people's children. Then a man from Belfast bought the
field and built houses in it --- not like their little brown houses but bright brick houses
with shining roofs. The children of the avenue used to play together in that field --- the
Devines, the Waters, the Dunns, little Keogh the \gl[a cripple]{cripple}{un infirme}, she and
her brothers and sisters. Ernest, however, never played: he was too grown up. Her father used
often to hunt them in out of the field with his \gl[a blackthorn stick]{blackthorn
stick}{une canne en prunellier}; but usually little Keogh used to keep \emph{nix} and call out
when he saw her father coming. Still they seemed to have been rather happy then. Her father
was not so bad then; and besides, her mother was alive. That was a long time ago; she and her
brothers and sisters were all grown up; her mother was dead. Tizzie Dunn was dead, too, and
the Waters had gone back to England. Everything changes. Now she was going to go away like the
others, to leave her home.

Home! She looked round the room, reviewing all its familiar objects which she had
\gl[to dust]{dusted}{épousseter} once a week for so many years, wondering where on earth all
the dust came from. Perhaps she would never see again those familiar objects from which she
had never dreamed of being divided. And yet during all those years she had never found out
the name of the priest whose yellowing photograph hung on the wall above the broken
\gl[a harmonium]{harmonium}{un harmonium} beside the coloured print of the promises made to
Blessed Margaret Mary Alacoque. He had been a school friend of her father. Whenever he showed
the photograph to a visitor her father used to pass it with a casual word:

``He is in Melbourne now.''

She had consented to go away, to leave her home. Was that wise? She tried to weigh each side
of the question. In her home anyway she had shelter and food; she had those whom she had known
all her life about her. Of course she had to work hard, both in the house and at business.
What would they say of her in the Stores when they found out that she had run away with a
fellow? Say she was a fool, perhaps; and her place would be filled up by advertisement. Miss
Gavan would be glad. She had always had an \gl[to have an edge on somebody]{edge}{avoir une
dent contre quelqu'un} on her, especially whenever there were people listening.

``Miss Hill, don't you see these ladies are waiting?''

``Look lively, Miss Hill, please.''

She would not cry many tears at leaving the Stores.
\end{texte}
\source{James Joyce, \og Eveline \fg, \emph{Dubliners}, Londres, Grant Richards, 1914.}

\partiel[15 min]{Questions}
\consigne{\en{Answer in English, in full sentences.}}

\begin{enumerate}[label=\arabic*.,itemsep=2mm,leftmargin=*]
  \item \en{What has become of the field where the children used to play?}
    \lignes{2}
    \corr{A man from Belfast bought it and built bright brick houses with shining roofs on
    it, unlike the little brown houses of the avenue.}
  \item \en{What has become of the children who played there?}
    \lignes{3}
    \corr{They are all grown up. Ernest is dead, Tizzie Dunn is dead too, the Waters have
    gone back to England, and Eveline is about to leave in her turn.}
  \item \en{``Still they seemed to have been rather happy then.'' What is the force of
    ``seemed'' and of ``rather''?}
    \lignes{3}
    \corr{Both weaken the sentence. She is not saying they were happy but that it looks that
    way from here, and only moderately so. Even her memory of childhood refuses to be simple,
    and the reason follows at once: her father was not so bad then.}
  \item \en{What does she find when she looks round the room?}
    \lignes{3}
    \corr{Familiar objects she has dusted once a week for years without ever wondering where
    the dust came from, and a yellowing photograph of a priest whose name she has never
    learned in all those years.}
  \item \en{List the reasons for staying and the reasons for leaving that she weighs.}
    \lignes{4}
    \corr{For staying: shelter, food, and people she has known all her life. For leaving:
    hard work at home and at the Stores, Miss Gavan's hostility, and the promise of respect
    as a married woman. What she never mentions in this passage is whether she loves Frank.}
  \item \en{Why does Joyce quote Miss Gavan's two sentences instead of describing her?}
    \lignes{3}
    \corr{Because the tone does the work: two orders, one of them humiliating in front of
    customers, tell us more about the daily life of the shop than a paragraph would. The
    whole story works by letting remembered phrases surface.}
\end{enumerate}

\note[Les tours du texte]{\emph{One time there used to be a field there}, \emph{they used to
play}, \emph{her father used often to hunt them}, \emph{used to keep nix} : l'habitude passée
de la fiche 2, cinq fois en dix lignes --- c'est ce qui donne au passage son air de mémoire
usée. Noter aussi \emph{used often to}, avec l'adverbe entre \emph{used} et \emph{to} : usage
d'époque, aujourd'hui \emph{often used to}.}

%% ===================================================================
\newpage
\section{Version}
\consigne{Traduisez la fin de la nouvelle. Eveline a rejoint Frank au port ; le bateau va
partir.}

\begin{version}
She stood among the \gl[to sway]{swaying}{ondoyer, osciller} crowd in the station at the North
Wall. He held her hand and she knew that he was speaking to her, saying something about the
\gl[the passage]{passage}{la traversée} over and over again. The station was full of soldiers
with brown \gl[baggages]{baggages}{les bagages} . Through the wide doors of the
\gl[a shed]{sheds}{un hangar} she caught a glimpse of the black mass of the boat, lying in
beside the quay wall, with \gl[illumined]{illumined}{éclairé} portholes. She answered nothing.
She felt her cheek pale and cold and, out of a \gl[a maze]{maze}{un labyrinthe} of distress,
she prayed to God to direct her, to show her what was her duty. The boat blew a long mournful
whistle into the mist. If she went, tomorrow she would be on the sea with Frank, steaming
towards Buenos Ayres. Their passage had been booked. Could she still draw back after all he
had done for her? Her distress awoke a \gl[nausea]{nausea}{une nausée} in her body and she
kept moving her lips in silent fervent prayer.

A bell \gl[to clang]{clanged}{résonner, retentir} upon her heart. She felt him seize her hand:

``Come!''

All the seas of the world tumbled about her heart. He was drawing her into them: he would
drown her. She gripped with both hands at the iron railing.

``Come!''

No! No! No! It was impossible. Her hands \gl[to clutch]{clutched}{agripper} the iron in
\gl[frenzy]{frenzy}{la frénésie}. Amid the seas she sent a cry of anguish!

``Eveline! Evvy!''

He rushed beyond the barrier and called to her to follow. He was shouted at to go on but he
still called to her. She set her white face to him, passive, like a helpless animal. Her eyes
gave him no sign of love or farewell or recognition.
\end{version}
\source{\emph{Ibid.}}

\lignes{22}

\corr{\textbf{Proposition de traduction.} Elle se tenait au milieu de la foule ondoyante, dans
la gare maritime de North Wall. Il lui tenait la main et elle savait qu'il lui parlait, qu'il
répétait sans cesse quelque chose au sujet de la traversée. La gare était pleine de soldats
avec des bagages bruns. Par les larges portes des hangars, elle aperçut un instant la masse
noire du bateau, amarré le long du quai, hublots éclairés. Elle ne répondit rien. Elle sentit
sa joue pâle et froide et, du fond d'un dédale de détresse, elle pria Dieu de la guider, de
lui montrer où était son devoir. Le bateau lança dans la brume un long coup de sirène
lugubre. Si elle partait, demain elle serait en mer avec Frank, filant vers Buenos Aires.
Leur passage était retenu. Pouvait-elle encore reculer, après tout ce qu'il avait fait pour
elle ? Sa détresse éveilla une nausée dans son corps, et elle continua de remuer les lèvres
en une prière fervente et muette.

Une cloche retentit sur son cœur. Elle le sentit lui saisir la main :

\og Viens ! \fg

Toutes les mers du monde déferlèrent autour de son cœur. Il l'entraînait dedans : il allait la
noyer. Elle s'agrippa des deux mains à la rambarde de fer.

\og Viens ! \fg

Non ! Non ! Non ! C'était impossible. Ses mains étreignirent le fer avec frénésie. Au milieu
des mers, elle poussa un cri d'angoisse !

\og Eveline ! Evvy ! \fg

Il se précipita au-delà de la barrière et lui cria de le suivre. On lui criait d'avancer, mais
il continuait de l'appeler. Elle tourna vers lui son visage blanc, passive, comme une bête
sans défense. Ses yeux ne lui donnèrent aucun signe d'amour, ni d'adieu, ni de
reconnaissance.}

\note[Choix de traduction 1]{\emph{He was shouted at to go on} : passif d'un verbe à
préposition --- impossible en français, où l'on revient à \og on \fg{} + actif (\og on lui
criait d'avancer \fg). C'est exactement l'exercice 1 de la fiche, en situation.}

\note[Choix de traduction 2]{\emph{she caught a glimpse of} : locution verbale figée,
\og elle entrevit \fg{} ou \og elle aperçut un instant \fg. Traduire \emph{glimpse} par
\og coup d'œil \fg{} inverserait le sens : c'est la vue qui est brève, pas le regard.}

\note[Choix de traduction 3]{\emph{to show her what was her duty} : interrogative indirecte
avec inversion, tour ancien pour \emph{what her duty was}. Le français dit \og où était son
devoir \fg.}

\note[Choix de traduction 4]{\emph{she kept moving her lips} : \emph{keep} + gérondif =
continuer à, ne pas cesser de (fiche 18). Ne pas traduire \emph{keep} par \og garder \fg.}

\note[Choix de traduction 5]{\emph{Her eyes gave him no sign of love or farewell or
recognition} : la négation porte sur le complément (\emph{no sign}), et les trois noms sont
coordonnés sans virgule finale. Le français répète la négation : \og ni d'adieu, ni de
reconnaissance \fg.}

%% ===================================================================
\partiel{Lexique à mémoriser}
\begin{itemize}[label=--,itemsep=0.8mm,leftmargin=*]\small
  \item \textbf{a footstep} — un bruit de pas \textit{(fiche 2)}
  \item \textbf{a pavement} — un trottoir \textit{(fiche 11)}
  \item \textbf{dust} — la poussière ; \textit{verbe} to dust \textit{« épousseter », adjectif} dusty
  \item \textbf{a shelter} — un abri ; \textit{verbe} to shelter
  \item \textbf{to weigh} — peser ; \textit{nom} weight\textit{, le} gh \textit{ne se prononce pas}
  \item \textbf{to consent} — consentir ; \textit{nom} consent
  \item \textbf{wise} — sage, avisé ; \textit{nom} wisdom
  \item \textbf{a fool} — un idiot ; \textit{adjectif} foolish \textit{(fiche 9)}
  \item \textbf{an advertisement} — une annonce, une publicité ; \textit{abrégé en} an ad
  \item \textbf{a crowd} — une foule \textit{(fiche 12)}
  \item \textbf{to sway} — osciller \textit{(fiche 8)}
  \item \textbf{a glimpse} — un aperçu ; to catch a glimpse of
  \item \textbf{mournful} — lugubre ; to mourn \textit{« porter le deuil »}
  \item \textbf{the mist} — la brume ; \textit{adjectif} misty
  \item \textbf{to draw back} — reculer ; drew, drawn
  \item \textbf{to drown} — noyer ; \textit{le} w \textit{ne se prononce pas}
  \item \textbf{a railing} — une rambarde, une grille
  \item \textbf{helpless} — sans défense ; help + -less
  \item \textbf{a farewell} — un adieu ; fare \textit{« aller » +} well
  \item \textbf{recognition} — la reconnaissance ; to recognise
\end{itemize}

\end{document}
